% `tex.web` §1215, §262 and §233: three things a `{restoring …}` line has to
% get right, all of them met at trip.tex line 213.
%
%     {\catcode`?=13 \font?xyzzy at0pt\font ? xyzzy scaled1?}
%
% §1257's `\font` opens with `get_r_token`, and §1215's `cur_cs` is nonzero for
% an *active character* too — they live in eqtb alongside the control
% sequences. So `\font?xyzzy` names a font `?` exactly as `\font\rip` names one
% `\rip`, and the trailing `?` on that line is a font identifier meaning
% \nullfont. We asked only for a control sequence, so the whole definition was
% dropped and `?` stayed undefined.
%
% §262 and §263 then print an active character *bare*:
%
%     else print(p-active_base)
%
% no escape, which is also how a document writes it. This engine keeps an
% active character's eqtb slot under a name beginning with a marker character,
% so `{restoring ?=undefined}` came out `{restoring \^^A?=undefined}` once the
% definition did happen — and every `\show~` with it.
%
% §233's region 4 is the third:
%
%     if n=cur_font_loc then print("current font")
%     …
%     print_char("="); print_esc(font_id_text(equiv(n)));
%
% The current font is an eqtb entry like any other, so a group restores it and
% §252 says so. We restored it — \fontname\font was right on either side — but
% off the save stack's trace, so every `{\bf …}` in a document was one
% `{restoring …}` line short of TeX's. `font_id_text` is the control sequence
% the font answers to, which is the *last* one bound to it, not the file it was
% read from.
\catcode`\{=1 \catcode`\}=2
\tracingonline=1 \tracingrestores=1 \errorcontextlines=5
\font\a=cmr10 \font\b=cmbx10 \a
\showthe\escapechar

% The plain form: a group that changes the font restores it, and says so.
{\b X}

% An active character as a font identifier, and one that cannot be loaded —
% so `?` means \nullfont, selecting it is a font selection like any other, and
% the group restores both it and the font.
{\catcode`?=13 \font?xyzzy at0pt\font ? xyzzy scaled1?}

% An active character as an ordinary definition, restored the same way. `!`
% rather than `~` on purpose: this engine's initial \catcode table is plain's
% and not §232's, so `~` starts out active here and *other* in TeX, and its
% restore could not agree before the initial-catcode fix. `!` is `other` in
% both engines.
{\catcode`!=13 \def!{bang}\show!}
\showthe\escapechar
\end
